% Generated draft: requires mathematical and source review.
\documentclass[11pt]{article}
\usepackage[margin=25mm]{geometry}
\usepackage{fontspec}
\setmainfont{Latin Modern Roman}
\usepackage{amsmath,amssymb,mathtools,unicode-math}
\setmathfont{Latin Modern Math}
\usepackage{enumitem,needspace,xurl,hyperref}
\setlength{\parindent}{0pt}
\setlength{\parskip}{5pt}
\setlength{\emergencystretch}{3em}
\newcommand{\N}{\mathbb{N}}
\newcommand{\Z}{\mathbb{Z}}
\newcommand{\Q}{\mathbb{Q}}
\newcommand{\R}{\mathbb{R}}
\newcommand{\C}{\mathbb{C}}
\newcommand{\F}{\mathbb{F}}
\newcommand{\A}{\mathbb{A}}
\newcommand{\PP}{\mathbb{P}}
\newcommand{\E}{\mathbb{E}}
\newcommand{\eps}{\varepsilon}
\newcommand{\dd}{\,\mathrm{d}}
\begin{document}
\section*{UnsolvedMath ID: 20000804}
\textbf{Source problem code:} catalogue entry
\subsection{Definitions and mathematical statement}
Fix a prime \(p\) and work over \(\overline{\F}_p\). On the Hilbert locus \(U\) of smooth geometrically integral nondegenerate curves \(C\subset\PP^r\), \(r\geq2\), of degree \(d\) and genus \(g\), the Gauss map
\[
\gamma_C:C\to\operatorname{Gr}(2,r+1)
\]
sends a point to its embedded tangent projective line. It is nonconstant on this locus and hence finite onto its image. Define
\[
Z_e=\{[C]\in U:[k(C):k(\gamma_C(C))]_{\mathrm{insep}}=p^e\},
\qquad e\geq0.
\]
Here the degree means the inseparable part, not an unstated assumption that the total degree is purely inseparable; \(e=0\) includes generically separable Gauss maps. These sets partition \(U\). The original says only smooth curves and prints positive \(e\) before taking a union over \(e\geq0\); degenerate, disconnected, or constant-Gauss-map cases require additional conventions outside this reconstructed locus.

The problem asks what geometric properties the \(Z_e\) have, whether this stratification can construct exotic Hilbert components existing only in characteristic \(p\), and how \(\operatorname{Gal}(\overline{\F}_p/\F_p)\) acts. It also suggests other stratifications, such as by the topological type of a ramification divisor. A precise use of exotic requires a chosen integral Hilbert model and a lifting/vertical-component criterion; the source does not specify one.

The attached aid reports that each fixed-degree set \(Z_e\) is constructible, defined over \(\F_p\), and Galois stable. It gives the necessary condition
\[
p^e\mid 2d+2g-2,
\]
and says each geometric Hilbert component meeting the indicated locus has one generic exponent. These are set-theoretic/constructibility statements, not claims that every \(Z_e\) is a closed subscheme or a Hilbert component.

For a concrete family take \(e\geq1\), \(q=p^e\), and an invertible three-by-three matrix \(B\), up to nonzero scalar. In \(\PP^2\), put
\[
F_B(X)=(X^{[q]})^TBX,
\qquad X^{[q]}=(X_0^q,X_1^q,X_2^q)^T.
\]
The aid reports that these smooth degree-\(q+1\) curves, of genus \(q(q-1)/2\), form a subvariety \(W_e\cong\operatorname{PGL}_3\) of dimension \(8\), closed relative to the smooth plane-curve Hilbert locus, contained in \(Z_e\). Their Gauss maps are exactly
\[
[X]\longmapsto[B^TX^{[q]}].
\]
The arithmetic-Frobenius convention \(a\mapsto a^p\) acts on parameters by \([B]\mapsto[B^{(p)}]\). This concerns embedded curves, not a quotient by projective equivalence. The restriction \(e\geq1\) is needed for this matrix parametrization; the separable stratum \(e=0\) is part of the general question but is not asserted to have the same description.

The aid emphasizes that \(W_e\) lies properly in the unique horizontal plane-curve Hilbert component. Consequently this explicit Frobenius family is not itself an exotic component. Constructibility, an inseparable Gauss map, or a characteristic-specific closed family alone does not prove existence of a new Hilbert component. These are supplied unreviewed partial results; the broader component-construction question remains unresolved by this family.
\subsection{Short English statement}
Study Hilbert-locus strata by the inseparable degree of a curve's Gauss map, their Galois action, and possible characteristic-specific components. The supplied aid gives constructibility and a Frobenius plane-curve family, while explicitly noting that the family is a proper subvariety of a horizontal component.
\subsection{Sources}
\begin{itemize}
\item Problem 3, \emph{Components of Hilbert Schemes}, recorded by Izzet Coskun and edited by Li Li after the AIM workshop of 19--23 July 2010, according to the supplied verification: \url{https://aimath.org/WWN/hilbertschemes/hilbertschemes.pdf}. The original remark also suggests ramification-divisor stratifications.
\item Local registry \texttt{data/unsolvedmath/problems.json}, permanent ID \(20000804\), exact source code \texttt{catalogue entry}; canonical \texttt{aim-arithmetic-geometry-notes.json}, notes[49].
\item The attached aid supplies the constructibility, divisibility, and matrix-family results without full external citations. Attempt 2 is explicitly machine-generated, not peer reviewed, and requires independent verification before citation.
\end{itemize}
\end{document}
